% Compile this document for the 20-minute talk. The paper entry point is main.tex.
% Spoken text is edited in the script .tex files, not in the rehearsal wrapper.
% Slides 1--4:   patches_talk_opening_script.tex
% Slides 5--9:   patches_talk_results_script.tex
% Slides 10--12: patches_talk_duality_script.tex
% Slides 13--17: patches_talk_patches_script.tex
% Slides 18--23: patches_talk_proof_script.tex
% Search each script file for the slide number or title. The .txt files are exports.
\ifdefined\TalkRehearsal
  % The rehearsal wrapper has already declared the document class.
\else
  \documentclass[t,aspectratio=169,11pt]{beamer}
\fi
\usetheme{Madrid}
\usecolortheme{default}
\usefonttheme{professionalfonts}
\usepackage[T1]{fontenc}
\usepackage[utf8]{inputenc}
\usepackage{lmodern,amsmath,amssymb,mathtools,booktabs,graphicx,pgfpages,tikz,adjustbox}
\usetikzlibrary{arrows.meta,calc,patterns.meta}
\definecolor{talkblue}{RGB}{31,65,133}
\definecolor{talkred}{RGB}{175,32,38}
\setbeamercolor{structure}{fg=talkblue}
\setbeamercolor{frametitle}{bg=talkblue,fg=white}
\setbeamercolor{title}{bg=talkblue,fg=white}
\setbeamercolor{block title}{bg=talkblue!12,fg=talkblue}
\setbeamercolor{block body}{bg=talkblue!4,fg=black}
\setbeamercolor{alerted text}{fg=talkred}
\setbeamertemplate{blocks}[rounded][shadow=false]
\setbeamertemplate{navigation symbols}{}
\setbeamercolor{itemize item}{fg=talkblue}
\setbeamertemplate{itemize item}{\tikz[baseline=-.55ex]\fill[talkblue] (0,0) circle (1.05pt);}
\setbeamertemplate{headline}{}
\setbeamersize{text margin left=7mm,text margin right=7mm}
\setbeamerfont{frametitle}{size=\large}
\setbeamerfont{block title}{size=\normalsize}
\setbeamerfont{footline}{size=\tiny}
\ifdefined\TalkRehearsal
  \setbeameroption{show notes on second screen=right}
  \setbeamertemplate{note page}{%
    \fontsize{10.5}{13.2}\selectfont
    \setlength{\parindent}{0pt}%
    \setlength{\parskip}{.65em}%
    \insertnote
  }
\else
  \setbeameroption{hide notes}
\fi
% Dark-red sentences in the rehearsal script are the lines to memorize.
\newcommand{\keyscript}[1]{{\color{talkred}#1}}
% Full spoken scripts for the opening, revised from the long introduction recording.
% Both the presentation and rehearsal versions read this file.

\newcommand{\OpeningTitleScript}{%
{\color{talkblue}\bfseries Slide 1: Opening}\par
{\small 0:00--0:25 (25 seconds)}\par\medskip
Hello everyone. This is joint work with Georg Menz, as part of a larger project where we ask the question: \keyscript{When does adding arbitrarily weak noise make a spin system ergodic?} Today I'll present one result from that project. To introduce the setting, let me begin with the Fredrickson--Andersen model, or FA-1f.%
}

\newcommand{\OpeningModelScript}{%
{\color{talkblue}\bfseries Slide 2: The FA-1f model}\par
{\small 0:25--1:25 (60 seconds)}\par\medskip
In FA-1f, each lattice site is either occupied, denoted by one, or vacant, denoted by zero.

Each site has a rate-one clock. \keyscript{When it rings, we replace the site's state by an independent Bernoulli-$p$ random variable, provided at least one neighbour is vacant.} The new state is one with probability p and zero with probability one minus p. If all neighbours are occupied, the update is blocked.

Here is the generator. I denote the constraint by $\lambda_i(\eta)$, and $E_{i,p}$ is the Bernoulli refresh operator at site i.

\keyscript{The Bernoulli-$p$ product measure is by construction invariant and reversible.}%
}

\newcommand{\OpeningQuestionScript}{%
{\color{talkblue}\bfseries Slide 3: Convergence to equilibrium}\par
{\small 1:25--2:20 (55 seconds)}\par\medskip
\keyscript{A long-standing conjecture says that FA-1f converges to this equilibrium from every initial configuration which contains at least one vacancy.} \keyscript{The fully occupied configuration is frozen.}

There are several partial results. For example, when p is less than one half, convergence is known from every Bernoulli product initial law whose occupation density is below one. The general statement for arbitrary initial configurations and all p remains open.

If true, the conjecture would leave only the fully occupied state as an obstruction to ergodicity.%
}

\newcommand{\OpeningNoiseScript}{%
{\color{talkblue}\bfseries Slide 4: Noisy FA-1f}\par
{\small 2:20--4:00 (100 seconds)}\par\medskip
\keyscript{Now add independent Bernoulli-$q$ refreshes at some small rate epsilon, without any constraint.} These destroy the fully occupied absorbing state.

\keyscript{If we believe the original conjecture, it is natural to expect that this perturbation makes FA-1f ergodic, however small epsilon is.} This is an example of the larger noise question.

One reason this is difficult is the lack of monotonicity. Monotonicity gives couplings, which are a standard tool for proving convergence. FA-1f is not monotone, so that standard argument is unavailable.

There is also a new difficulty: if q differs from p, the noise no longer preserves the original equilibrium, and the invariant measure is generally neither known explicitly nor reversible.

Our contribution is a way to prove ergodicity despite these obstacles, for a more general class of spin systems. To state our results, let me first introduce the notation.%
}

% Full spoken results scripts, revised from the second results recording.
% The presentation and rehearsal versions read this file.

\newcommand{\ResultsSettingScript}{%
{\color{talkblue}\bfseries Slide 5: General setting}\par
{\small 4:00--4:25 (25 seconds)}\par\medskip
We write a general spin system in the refresh form borrowed from kinetically constrained models. \keyscript{At site i, $\lambda_i(\eta)$ is the update rate, and $\rho_i(\eta)$ is the probability that the new spin is one.}

\keyscript{Every spin system can be written this way.} This representation is convenient for our graphical construction.%
}

\newcommand{\ResultsPositivityScript}{%
{\color{talkblue}\bfseries Slide 6: Patch positivity and the density profile}\par
{\small 4:25--5:35 (70 seconds)}\par\medskip
To describe the spin systems for which our results apply, we use monomial coefficients of the update rate and refresh parameter; see the hat notation on the slide. \keyscript{We call the spin systems defined by these two systems of inequalities patch-positive spin systems.} This is an explicit condition on the rates; its probabilistic meaning will be derived later from our graphical construction of patches.

The class includes FA-1f and related models such as the biased annihilating branching process and the contact process.

\keyscript{Each patch-positive spin system also has an explicit density profile, shown on the slide.} This formula too has a natural meaning in our graphical construction of patches. For FA-1f, this density is simply p.%
}

\newcommand{\ResultsLawsScript}{%
{\color{talkblue}\bfseries Slide 7: The initial laws}\par
{\small 5:35--6:15 (40 seconds)}\par\medskip
\keyscript{In the main result, we prove convergence to a common invariant limit from a class of initial laws defined using $\boldsymbol p^\star$.} To define this class, we introduce monomials centered at $\boldsymbol p^\star$; see the formula on the slide. \keyscript{Let $\mathcal M_\star$ consist of the laws with nonnegative centered moments on every finite set.} In particular, $\mathcal M_\star$ contains every Bernoulli product law with density at least $\boldsymbol p^\star$.

We can cheaply extend our results to a larger class, $\mathcal M_-$. It contains the laws that enter $\mathcal M_\star$ when averaged with $\delta_{\mathbf1}$.%
}

\newcommand{\ResultsTheoremScript}{%
{\color{talkblue}\bfseries Slide 8: Convergence with a pure-death component}\par
{\small 6:15--7:15 (60 seconds)}\par\medskip
Here is the main result. \keyscript{Suppose the system is patch-positive, its neighbourhood graph has polynomial growth, and it contains independent one-to-zero transitions at a positive rate epsilon.} \keyscript{Then every initial law in $\mathcal M_-$ converges exponentially to the same invariant measure.} The estimate is uniform on $\mathcal M_{-,K}$.

To obtain ergodicity, we need convergence from every initial law. This follows under either of two additional conditions. In the high-temperature or high-vacancy regime, meaning $p_i^\star\leq1/2$, $\mathcal M_{-,1}$ contains every probability law, so we get uniform exponential ergodicity.

Alternatively, we can upgrade the result to uniform ergodicity, with no condition on $\boldsymbol p^\star$, but assuming the neighbourhood graph is orientable. In both cases, epsilon may be arbitrarily small.%
}

\newcommand{\ResultsApplicationsScript}{%
{\color{talkblue}\bfseries Slide 9: Noisy FA-1f: two ergodicity regimes}\par
{\small 7:15--8:00 (45 seconds)}\par\medskip
Let's return to our noisy FA-1f question. The noisy system is patch-positive exactly when q is at least p. With q below one, the noise gives pure one-to-zero transitions.

\keyscript{This proves our noisy ergodicity claim in two cases: the high-vacancy regime p at most one half, and when the neighbourhood graph is orientable.}

\keyscript{The noise can be arbitrarily weak, provided p is at most q and q is below one.} Let me now explain the proof.%
}

% Spoken duality scripts, based on the standard-duality recording.
% Finite-volume/localization notation is suppressed in the talk.

\newcommand{\DualityProcessScript}{%
{\color{talkblue}\bfseries Slide 10: Duality: tracing dependence backwards}\par
{\small 8:00--9:25 (85 seconds)}\par\medskip
Our proof starts from duality because it is a standard way to study convergence without coupling and without a priori knowing the limiting measure. Classically, this proves ergodicity when interactions are sufficiently weak compared with independent updates.

\keyscript{The dual consists of a finite set $A_s$ and a sign $\sigma_s$.} Starting from $A$, think of $A_s$ as tracing backwards which earlier spins can influence the spins in $A$ at time t.

The three types of updates are deaths, splits and births. The picture shows an example trajectory starting from sites three and four. A death removes site four. A split removes site three and adds sites two and four. A birth then adds sites one and three, while leaving sites two and four active.

\keyscript{Duality gives the Feynman--Kac representation below.} $\chi_A$ is the product of the spins in $A$. The representation also features a sign process $\sigma_t$ and an exponential weight involving a potential $V$.%
}

\newcommand{\DualityClassicalScript}{%
{\color{talkblue}\bfseries Slide 11: The usual duality argument}\par
{\small 9:25--10:25 (60 seconds)}\par\medskip
For ergodicity, we want $P_t\chi_A$ to converge to a constant with respect to $\eta$, for every finite $A$.

\keyscript{Duality lets us split $P_t\chi_A$ into a term which does not depend on the initial configuration and a term which is negligible if the process is ergodic.} We make this split according to whether $A_t$ is empty.

The classical argument assumes that the potential is bounded above by a strictly negative constant while the dual set is nonempty. On that event, the exponential weight is at most $e^{-\gamma t}$. The second term vanishes, giving a limit independent of the initial configuration.%
}

\newcommand{\DualityObstructionScript}{%
{\color{talkblue}\bfseries Slide 12: Why the standard estimate fails}\par
{\small 10:25--11:00 (35 seconds)}\par\medskip
For weakly noisy FA-1f, this estimate fails. \keyscript{The potential is positive rather than negative, so the exponential weight can be very large.}

\keyscript{To show that the term depending on $\eta$ is negligible, we have to control the cancellations from $\sigma_t$.} The large weight makes that control particularly demanding because even rare events can make substantial contributions.

Next we consider how to average over part of the graphical construction to capture these cancellations.%
}

% Spoken scripts for the patch-representation section, based on the recording.

\newcommand{\PatchesConditioningScript}{%
{\color{talkblue}\bfseries Slide 13: How much of the dual should we reveal?}\par
{\small 11:00--11:50 (50 seconds)}\par\medskip
First focus on the signed exponential weight. \keyscript{We want a sigma-algebra $\mathcal G_t$ with two properties: the conditional mean should be nonnegative, and it should be explicit enough to analyse.}

If we reveal nothing, we average all the cancellations. The mean is nonnegative: it is $P_t\chi_A$ evaluated at the all-one configuration. But we have gained no way to compute it.

If we reveal the entire dual path, the weight is known, but it still carries its random sign. \keyscript{We need an intermediate amount of information.}%
}

\newcommand{\PatchesGeometryScript}{%
{\color{talkblue}\bfseries Slide 14: The interactions define the patches}\par
{\small 11:50--12:50 (60 seconds)}\par\medskip
Here is our choice. In the graphical construction, births and splits are interactions between sites; deaths are not. \keyscript{Let $\mathcal G_t$ be the sigma-algebra generated by the spacetime positions of the interactions and the sites involved, up to time t.} \keyscript{It does not reveal whether an interaction is a birth or a split, or any deaths.}

The left picture shows the dual realization. On each site line, the incoming and outgoing interactions divide time into intervals. These are the patches shown on the right.

There are four bulk types, according to whether the interaction at each end is incoming or outgoing. Patches reaching the observation time are cut there and become end patches. The thick red line separates bulk patches from end patches.%
}

\newcommand{\PatchesIndependenceScript}{%
{\color{talkblue}\bfseries Slide 15: Conditional independence across patches}\par
{\small 12:50--13:35 (45 seconds)}\par\medskip
The sigma-algebra $\mathcal G_t$ has the properties we wanted. \keyscript{First, it provides enough information for explicit computation: conditional on $\mathcal G_t$, the law of the dual factorizes over patches.} The remaining randomness in different patches is independent.

\keyscript{Each patch has an explicit one-site law $\pP_P^{\mathrm{con}}$, which is why $\mathcal G_t$ satisfies the second condition.}%
}

\newcommand{\PatchesInteriorsScript}{%
{\color{talkblue}\bfseries Slide 16: The four bulk patch types}\par
{\small 13:35--13:45 (10 seconds)}\par\medskip
This figure shows what the process can do inside each patch type. \keyscript{The possibilities are simple enough to compute each patch law explicitly.}%
}

\newcommand{\PatchesRepresentationScript}{%
{\color{talkblue}\bfseries Slide 17: The patch representation}\par
{\small 13:45--15:00 (75 seconds)}\par\medskip
The signed Feynman--Kac weight is itself a product over patches. \keyscript{Conditional independence with respect to $\mathcal G_t$ lets us average each factor separately.} A bulk patch contributes a number $C(P)$. An end patch contributes $C(z,P)$, where z is the initial spin at that site in the original system.

We call the product of these contributions $W_t^\mu$. The initial law $\mu$ acts on the whole product of end contributions, since the initial spins may be correlated.

\keyscript{We can give all these factors explicitly, and all dependence on the initial law remains in the end patches.} This is the patch representation. Next we use patch positivity to control the signs of these contributions.%
}

% Full spoken proof scripts, revised from the second proof-end recording.
% The presentation and rehearsal versions read this file.

\newcommand{\ProofPositivityScript}{%
{\color{talkblue}\bfseries Slide 18: Patch positivity and the meaning of p-star}\par
{\small 15:00--15:35 (35 seconds)}\par\medskip
Here is the meaning of the two conditions we stated earlier. \keyscript{Patch positivity is the condition that $C(P)$ is nonnegative for every bulk patch P.} \keyscript{The profile $\boldsymbol p^\star$ gives the smallest densities at which the end contributions are nonnegative.}

It defines the class $\mathcal M_\star$, for which the averaged product of end contributions is also nonnegative. All the factors in the patch factorization are nonnegative, so the sigma-algebra $\mathcal G_t$ captures all the cancellations we wanted.%
}

\newcommand{\ProofSplitScript}{%
{\color{talkblue}\bfseries Slide 19: Splitting by late interactions}\par
{\small 15:35--16:30 (55 seconds)}\par\medskip
The proof now follows the structure of the classical duality argument. We split $(\mu P_t)(\chi_A)$ into two terms: one with weak dependence on the initial law, and one whose contribution is negligible.

The extinction event $\{A_t=\varnothing\}$ is no longer the right event to split by: it is not $\mathcal G_t$-measurable, because the process A is hidden behind the patches. But we have a good proxy that is $\mathcal G_t$-measurable: $L_t$, the event that there is no interaction between t over two and t.

\keyscript{On $L_t$, we show weak dependence on $\mu$.} \keyscript{On its complement, we show that the whole contribution is small.}%
}

\newcommand{\ProofLateScript}{%
{\color{talkblue}\bfseries Slide 20: Late interaction term}\par
{\small 16:30--18:00 (90 seconds)}\par\medskip
To control the late-interaction term, we make the following observation. \keyscript{A spin system and its less noisy version, obtained by removing the independent one-to-zero transitions at rate epsilon, have exactly the same dual process.} The less noisy system has a larger Feynman--Kac potential.

They therefore share the random patches, but the less noisy system has larger patch contributions. We use primes for that system. For a patch ending with an outgoing interaction, the ratio is exactly $e^{-\varepsilon\abs P}$.

A late interaction at $u>t/2$ has an ancestry trail of patches ending with outgoing interactions, from time zero to u. \keyscript{Multiplying the ratios along the trail gives $e^{-\varepsilon u}\leq e^{-\varepsilon t/2}$.}

This gives the displayed chain of inequalities. The primed expectation, after dropping the event, is $(\mu P_t')(\chi_A)$, where $P_t'$ is the semigroup of the less noisy system. It is at most one, because $\chi_A$ is bounded by one. The late-interaction term is therefore negligible.%
}

\newcommand{\ProofEndScript}{%
{\color{talkblue}\bfseries Slide 21: No-late-interaction term}\par
{\small 18:00--18:45 (45 seconds)}\par\medskip
\keyscript{On $L_t$, there is no interaction after t over two, so every end patch extends for at least t over two.} These are the only patches that carry dependence on the initial law.

\keyscript{In the contribution from this event, the dependence on the spin z in $C(z,P)$ decays exponentially with the added patch length.} All the end patches are long, so we can compare the contributions from two initial laws and obtain the displayed estimate.

The errors accumulate over the number of end patches, $\abs{\mathcal E_t}$. This factor is controlled using polynomial growth of the lattice. This is where that assumption enters.%
}

\newcommand{\ProofLimitScript}{%
{\color{talkblue}\bfseries Slide 22: The common invariant limit}\par
{\small 18:45--19:20 (35 seconds)}\par\medskip
\keyscript{The two estimates give exponential loss of dependence on the initial law and hence convergence from $\mathcal M_\star$ to the same invariant law $\pi$.} We identify the common limit by comparing with an invariant law in $\mathcal M_\star$. Such a law exists because $\mathcal M_\star$ is preserved by the spin-system semigroup.

\keyscript{A simple linearity trick extends the convergence from $\mathcal M_\star$ to the larger class $\mathcal M_-$.} This proves the theorem.%
}

\newcommand{\ProofErgodicityScript}{%
{\color{talkblue}\bfseries Slide 23: From convergence to ergodicity}\par
{\small 19:20--20:00 (40 seconds)}\par\medskip
Finally, the theorem gives full ergodicity in the two regimes stated earlier. \keyscript{The first case, $p_i^\star\leq1/2$ for every i, is immediate because $\mathcal M_{-,1}$ contains every probability law.} We get uniform exponential ergodicity from the theorem.

\keyscript{On an orientable neighbourhood graph, a separate burn-in argument shows that every initial law approaches $\mathcal M_\star$, in the sense of the displayed bound on centered moments.} We already know convergence to the common invariant law from that class, which gives uniform ergodicity. Thank you.%
}

\newcommand{\MainFrames}{23}
\setbeamertemplate{footline}{%
  \leavevmode\hbox{%
    \begin{beamercolorbox}[wd=.29\paperwidth,ht=2.6ex,dp=1.1ex,leftskip=2mm]{author in head/foot}M. G{\l}uchowski\end{beamercolorbox}%
    \begin{beamercolorbox}[wd=.54\paperwidth,ht=2.6ex,dp=1.1ex,center]{title in head/foot}Patch positivity and ergodicity\end{beamercolorbox}%
    \begin{beamercolorbox}[wd=.17\paperwidth,ht=2.6ex,dp=1.1ex,center]{date in head/foot}\insertframenumber\,/\,\MainFrames\end{beamercolorbox}}}
\newcommand{\Z}{\mathbb Z}
\newcommand{\E}{\mathbb E}
\newcommand{\pP}{\mathbb P}
\newcommand{\ind}{\mathbf1}
\newcommand{\one}{\mathbf1}
\newcommand{\Lcal}{\mathcal L}
\newcommand{\Dcal}{\mathcal D}
\newcommand{\Mstar}{\mathcal M_\star}
\newcommand{\Mminus}{\mathcal M_-}
\newcommand{\bb}[1]{\left(#1\right)}
\newcommand{\Cb}[1]{\left[#1\right]}
\newcommand{\cb}[1]{\left\{#1\right\}}
\newcommand{\abs}[1]{\left|#1\right|}
\newcommand{\norm}[1]{\left\Vert#1\right\Vert}
\newcommand{\smallline}[1]{\par\smallskip{\small #1\par}}
\newenvironment{talkpoints}{\begin{itemize}\setlength{\itemsep}{.65em}}{\end{itemize}}
\setlength{\abovedisplayskip}{.65em}
\setlength{\belowdisplayskip}{.55em}
\newcommand{\talkfigure}[3]{\begingroup\rmfamily\adjustbox{width=#2,max height=#3}{\input{talk_figures/#1.tikz}}\endgroup}
\newcommand{\paperfigure}[2]{\talkfigure{#1}{\linewidth}{#2}}
\title[Patch positivity and ergodicity]{Patch positivity and applications\\to kinetically constrained spin models}
\author{Maciej G{\l}uchowski}
\institute{University of California, Los Angeles\\[2mm]Joint work with Georg Menz}
\date{AMS Special Session on Lattice Models of Statistical Mechanics\\Tempe, Arizona \quad November 2026}
\hypersetup{pdftitle={Patch positivity and applications to kinetically constrained spin models},pdfauthor={Maciej Gluchowski},pdfsubject={Draft 20-minute talk; joint work with Georg Menz}}
\begin{document}

% 1. 0:00--0:25
\begin{frame}[c,label=opening-title]
\titlepage
% Spoken text: patches_talk_opening_script.tex; search for OpeningTitleScript.
\note{\OpeningTitleScript}
\end{frame}

% 2. 0:25--1:25
\begin{frame}[label=opening-model]{The Fredrickson--Andersen model (FA-1f)}
\begin{talkpoints}
\item $\eta\in\{0,1\}^{\Z^d}$, with $0$ vacant and $1$ occupied.
\item At rate one, site $i$ refreshes to $\operatorname{Bern}(p)$
\textbf{if at least one neighbour is vacant.}
\begin{equation*}
1\text{ with probability }p,\qquad 0\text{ with probability }1-p.
\end{equation*}
\end{talkpoints}
\medskip
\begin{equation*}
\Lcal_{\mathrm{FA}}f
=\sum_i\underbrace{\left(1-\prod_{j\sim i}\eta(j)\right)}_{\lambda_i(\eta)}
\bigl(E_{i,p}-I\bigr)f.
\end{equation*}
\smallline{$\lambda_i(\eta)$ is the constraint; $E_{i,p}$ is the Bernoulli refresh operator at $i$.}
\medskip
\begin{equation*}
\mu_p=\bigotimes_i\operatorname{Bern}(p)\quad\text{is invariant and reversible},
\qquad 0<p<1.
\end{equation*}
% Spoken text: patches_talk_opening_script.tex; search for OpeningModelScript.
\note{\OpeningModelScript}
\end{frame}

% 3. 1:25--2:20
\begin{frame}[label=opening-question]{Convergence to equilibrium}
\begin{block}{The FA-1f convergence conjecture}
For every initial configuration which contains at least one vacancy,
\begin{equation*}
\delta_\eta P_t\Longrightarrow\mu_p,\qquad \eta\neq\one.
\end{equation*}
\end{block}
\medskip
\begin{talkpoints}
\item The fully occupied configuration is frozen:
\begin{equation*}
\delta_{\one}P_t=\delta_{\one}.
\end{equation*}
\item A partial result in the high-vacancy regime:
\begin{equation*}
0<p<\tfrac12,\quad 0\leq u<1
\quad\Longrightarrow\quad \mu_uP_t\Longrightarrow\mu_p.
\end{equation*}
\smallline{Blondel et al. (2013), Theorem 2.1 and Remark 2.3.}
\end{talkpoints}
% Spoken text: patches_talk_opening_script.tex; search for OpeningQuestionScript.
\note{\OpeningQuestionScript}
\end{frame}

% 4. 2:20--4:00
\begin{frame}[label=opening-noise]{Noisy FA-1f}
\begin{equation*}
\Lcal^{\varepsilon,q}
=\Lcal_{\mathrm{FA}}+\varepsilon\sum_i(E_{i,q}-I),
\qquad \varepsilon>0,\quad 0\leq q<1.
\end{equation*}
\begin{talkpoints}
\item Independent refreshes create vacancies even in the fully occupied state.
\end{talkpoints}
\medskip
\begin{block}{The noisy ergodicity question}
Is there a common invariant law $\pi$ such that
\begin{equation*}
\delta_\eta P_t^{\varepsilon,q}\Longrightarrow\pi
\qquad\text{for every }\eta?
\end{equation*}
\end{block}
% Spoken text: patches_talk_opening_script.tex; search for OpeningNoiseScript.
\note{\OpeningNoiseScript}
\end{frame}

% 5. 4:00--4:25
\begin{frame}[label=results-setting]{General setting}
\begin{equation*}
\eta\in\{0,1\}^{\Lambda},\qquad
\Lcal f=\sum_{i\in\Lambda}\lambda_i(\eta)
\bigl(E_{i,\rho_i(\eta)}-I\bigr)f.
\end{equation*}
\medskip
\begin{talkpoints}
\item $\lambda_i(\eta)\geq0$: update rate at site $i$.
\item $\rho_i(\eta)\in[0,1]$: probability of setting the new spin to $1$.
\end{talkpoints}
\medskip
Every spin system admits this refresh representation.
% Spoken text: patches_talk_results_script.tex; search for ResultsSettingScript.
\note{\ResultsSettingScript}
\end{frame}

% 6. 4:25--5:35
\begin{frame}[label=results-positivity]{Patch positivity and the density profile}
\smallline{Monomial coefficients of a function of the spins in $N(i)$:}
\begin{equation*}
g(\eta)=\sum_{S\subseteq N(i)}\widehat g(S)\prod_{j\in S}\eta(j).
\end{equation*}
\begin{block}{Patch positivity: an explicit condition on the rates}
\begin{equation*}
\widehat\lambda_i(S)\leq0,\qquad
\widehat{\lambda_i\bigl(\rho_i-\rho_i(\mathbf0)\bigr)}(S)\geq0,
\quad \varnothing\neq S\subseteq N(i).
\end{equation*}
\end{block}
\smallline{Examples: FA-1f, biased annihilating branching process, contact process.}
\medskip
Each patch-positive system has an explicit density profile $\boldsymbol p^\star\in[0,1]^\Lambda$:
\begin{equation*}
p_i^\star=\max\!\left(\{0\}\cup
\left\{\frac{\widehat{\lambda_i\rho_i}(S)}{\widehat\lambda_i(S)}:
\substack{\varnothing\neq S\subseteq N(i)\\\widehat\lambda_i(S)<0}\right\}\right).
\end{equation*}
\smallline{For FA-1f, $p_i^\star=p$.}
% Spoken text: patches_talk_results_script.tex; search for ResultsPositivityScript.
\note{\ResultsPositivityScript}
\end{frame}

% 7. 5:35--6:15
\begin{frame}[label=results-laws]{Some extra notation}
\begin{talkpoints}
\item $\displaystyle\chi_B^\star(\eta)=\prod_{i\in B}\bigl(\eta(i)-p_i^\star\bigr),
\qquad B\subset\Lambda\text{ finite}.$
\item $\displaystyle\Mstar=\cb{\mu:\ \mu(\chi_B^\star)\geq0\text{ for every finite }B}.$
\par\medskip
$\Mstar$ contains every product law with density at least $\boldsymbol p^\star$.
\item $\displaystyle\mathcal M_{-,K}=\cb{\mu:\ \frac{\mu+K\delta_{\one}}{1+K}\in\Mstar},
\qquad K\geq0.$
\begin{equation*}
\Mminus=\bigcup_{K\geq0}\mathcal M_{-,K}.
\end{equation*}
\end{talkpoints}
% Spoken text: patches_talk_results_script.tex; search for ResultsLawsScript.
\note{\ResultsLawsScript}
\end{frame}

% 8. 6:15--7:15
\begin{frame}[label=results-theorem]{Main result}
\begin{block}{Theorem (G{\l}uchowski--Menz)}
Assume patch positivity, polynomial graph growth with exponent $D$, and independent $1\to0$ transitions at rate $\varepsilon>0$.
\smallskip
There is an invariant law $\pi$ such that, for every local $f$,
\begin{equation*}
\sup_{\mu\in\mathcal M_{-,K}}
\abs{(\mu P_t)(f)-\pi(f)}
\leq (1+2K)C_f(1+t)^D e^{-\varepsilon t/2}.
\end{equation*}
\end{block}
\medskip
\textbf{Ergodicity under either additional condition:}
\par\medskip
\begin{tabular}{@{}p{.44\textwidth}@{\quad$\Longrightarrow$\quad}p{.45\textwidth}@{}}
High temperature / high vacancy:\newline $p_i^\star\leq\tfrac12$ for every $i$
& uniform exponential ergodicity\\[1.5mm]\midrule\addlinespace[1.5mm]
Orientable neighbourhood graph
& uniform ergodicity for every $\boldsymbol p^\star$
\end{tabular}
\par\medskip
\smallline{In the first case, $\mathcal M_{-,1}$ contains every initial law. The rate $\varepsilon$ may be arbitrarily small.}
% Spoken text: patches_talk_results_script.tex; search for ResultsTheoremScript.
\note{\ResultsTheoremScript}
\end{frame}

% 9. 7:15--8:00
\begin{frame}[label=results-applications]{Results for noisy FA-1f}
\begin{equation*}
\Lcal^{\varepsilon,q}=\Lcal_{\mathrm{FA}}+\varepsilon\sum_i(E_{i,q}-I),
\qquad \varepsilon>0.
\end{equation*}
\begin{equation*}
\text{the noisy system is patch-positive}\quad\Longleftrightarrow\quad q\geq p.
\end{equation*}
\smallline{For $q<1$, the noise gives pure $1\to0$ transitions.}
\medskip
\textbf{With $p\leq q<1$, the noisy system is ergodic in two cases:}
\par\medskip
\begin{tabular}{@{}p{.44\textwidth}@{\quad$\Longrightarrow$\quad}p{.45\textwidth}@{}}
High temperature / high vacancy:\newline $p\leq\tfrac12$
& uniform exponential ergodicity\\[1.5mm]\midrule\addlinespace[1.5mm]
Orientable neighbourhood graph\newline (in particular, East)
& uniform ergodicity for every $p\in(0,1)$
\end{tabular}
\par\medskip
\smallline{The noise strength $\varepsilon$ may be arbitrarily small.}
% Spoken text: patches_talk_results_script.tex; search for ResultsApplicationsScript.
\note{\ResultsApplicationsScript}
\end{frame}

% 10. 8:00--9:25
\begin{frame}[label=duality-process]{Duality}
Finite set $A_s\subset\Lambda$ and sign $\sigma_s\in\{-1,1\}$.
\begin{equation*}
\chi_A(\eta)=\prod_{i\in A}\eta(i),
\qquad (A_0,\sigma_0)=(A,1).
\end{equation*}
\begin{center}
\talkfigure{dual-idea}{.94\textwidth}{3.45cm}
\end{center}
% As in the paper, first read the signed identity in finite volume.
% Localization and passage to infinite volume are suppressed in this talk.
\begin{equation*}
P_t\chi_A(\eta)
=\E_A\!\left[
\sigma_t e^{\int_0^t V(A_u)\,du}\chi_{A_t}(\eta)
\right].
\end{equation*}
% Spoken text: patches_talk_duality_script.tex; search for DualityProcessScript.
\note{\DualityProcessScript}
\end{frame}

% 11. 9:25--10:25
\begin{frame}[label=duality-classical]{The usual duality argument...}
\begin{align*}
P_t\chi_A(\eta)
&=\underbrace{\E_A\!\left[
\sigma_t e^{\int_0^t V(A_u)\,du}\ind_{\{A_t=\varnothing\}}
\right]}_{\text{independent of }\eta\text{, since }\chi_\varnothing=1}
\\[3mm]
&\quad+\E_A\!\left[
\sigma_t e^{\int_0^t V(A_u)\,du}\chi_{A_t}(\eta)
\ind_{\{A_t\neq\varnothing\}}
\right].
\end{align*}
\medskip
\begin{talkpoints}
\item If $V(B)\leq-\gamma\ind_{\{B\neq\varnothing\}}$ for some $\gamma>0$, then
\begin{equation*}
e^{\int_0^tV(A_u)\,du}\ind_{\{A_t\neq\varnothing\}}
\leq e^{-\gamma t}.
\end{equation*}
\item The term depending on $\eta$ vanishes exponentially.
\end{talkpoints}
% Spoken text: patches_talk_duality_script.tex; search for DualityClassicalScript.
\note{\DualityClassicalScript}
\end{frame}

% 12. 10:25--11:00
\begin{frame}[label=duality-obstruction]{...fails}
For noisy FA-1f,
\begin{equation*}
V(A)=\bigl(2p-\varepsilon(1-q)\bigr)\abs A>0
\end{equation*}
when $A\neq\varnothing$ and $\varepsilon(1-q)<2p$.
\par\bigskip
The term still depending on $\eta$ is
\begin{equation*}
\E_A\!\left[
\sigma_t e^{\int_0^t V(A_u)\,du}\chi_{A_t}(\eta)
\ind_{\{A_t\neq\varnothing\}}
\right].
\end{equation*}
\medskip
\begin{talkpoints}
\item We need control of the cancellations carried by $\sigma_t$.
\item Even rare events can make substantial contributions.
\end{talkpoints}
% Spoken text: patches_talk_duality_script.tex; search for DualityObstructionScript.
\note{\DualityObstructionScript}
\end{frame}

% 13. 11:00--11:50
\begin{frame}[label=patches-conditioning]{How to capture cancellations?}
Choose $\mathcal G_t$ so that the conditional expectation
\begin{equation*}
\E_A\!\left[\sigma_t e^{\int_0^t V(A_u)\,du}\,\middle|\,\mathcal G_t\right]
\end{equation*}
has both properties:
\medskip
\begin{columns}[T,onlytextwidth]
\begin{column}{.47\textwidth}
\begin{block}{$\mathcal{G}_t$ small}
A nonnegative conditional mean.
\end{block}
\end{column}
\begin{column}{.47\textwidth}
\begin{block}{$\mathcal{G}_t$ rich}
An explicit conditional mean.
\end{block}
\end{column}
\end{columns}
\bigskip
\begin{talkpoints}
\item Reveal nothing,~$\mathcal{G}_t = \cb{\varnothing,\Omega}$: nonnegative mean $P_t\chi_A(\one)$, hard to analyse.
\item Reveal everything,~$\mathcal{G}_t=\mathcal{F}_t$: trivially explicit but signed.
\end{talkpoints}
% Spoken text: patches_talk_patches_script.tex; search for PatchesConditioningScript.
\note{\PatchesConditioningScript}
\end{frame}

% 14. 11:50--12:50
\begin{frame}[label=patches-geometry]{Patches}
\begin{equation*}
\mathcal G_t=\sigma\!\left(\text{spacetime positions and sites of interactions up to }t\right).
\end{equation*}
\smallline{The birth-or-split labels and deaths are not revealed.}
\begin{center}
\talkfigure{patch-geometry}{.96\textwidth}{5.25cm}
\end{center}
% Spoken text: patches_talk_patches_script.tex; search for PatchesGeometryScript.
\note{\PatchesGeometryScript}
\end{frame}

% 15. 12:50--13:35
\begin{frame}[label=patches-independence]{Conditional independence across patches}
\begin{talkpoints}
    \item The process, conditional on $\mathcal G_t$, is independent in each patch, meaning the process law factorizes
\end{talkpoints}
\bigskip
\begin{equation*}
\operatorname{Law}\!\left((A_s,\sigma_s)_{0\leq s\leq t}\,\middle|\,\mathcal G_t\right)
=\bigotimes_{\text{patches }P}\pP_P^{\mathrm{con}}.
\end{equation*}
\bigskip
Each patch has an explicit one-site law $\pP_P^{\mathrm{con}}$.
% Spoken text: patches_talk_patches_script.tex; search for PatchesIndependenceScript.
\note{\PatchesIndependenceScript}
\end{frame}

% 16. 13:35--13:45
\begin{frame}[label=patches-interiors]{Patch law~$\pP_P^{\mathrm{con}}$}
\begin{center}
\talkfigure{patch-interiors-all}{.98\textwidth}{6.65cm}
\end{center}
% Spoken text: patches_talk_patches_script.tex; search for PatchesInteriorsScript.
\note{\PatchesInteriorsScript}
\end{frame}

% 17. 13:45--15:00
\begin{frame}[label=patches-representation]{The patch representation}
We define
\begin{equation*}
\boxed{\displaystyle
W_t^\mu
=\prod_{P\in\mathcal B_t}C(P)\;
\mu\!\left[\prod_{P\in\mathcal E_t}C\bigl(\eta(i(P)),P\bigr)\right].}
\end{equation*}
\begin{equation*}
(\mu P_t)(\chi_A)=\E_A[W_t^\mu].
\end{equation*}
\bigskip
\begin{tabular}{@{}ll@{}}
$\mathcal B_t$ & bulk patches, ending before $t$\\[2mm]
$\mathcal E_t$ & end patches, reaching $t$\\[2mm]
$C(P),\ C(z,P)$ & contributions to~$\E_A\Cb{\sigma_t e^{\int_0^t V(A_s)ds}\chi_{A_t}\mid\mathcal{G}_t}$ from bulk\\
& and end patches
\end{tabular}
\bigskip
\begin{center}
Only the end contributions depend on initial configuration.
\end{center}
% Spoken text: patches_talk_patches_script.tex; search for PatchesRepresentationScript.
\note{\PatchesRepresentationScript}
\end{frame}

% 18. 15:00--15:35
\begin{frame}[label=proof-positivity]{\texorpdfstring{Patch positivity and $p^\star$}{Patch positivity and the density profile}}
\begin{talkpoints}
\item $\displaystyle\text{patch positivity}\quad\Longleftrightarrow\quad C(P)\geq0
\quad\text{for every bulk patch }P.$
\item $\displaystyle p_i^\star
=\min\cb{z\in[0,1]:C(z,P)\geq0\text{ for every end patch at }i}.$
\item $\displaystyle\mu\in\Mstar\quad\Longrightarrow\quad
\mu\!\left[\prod_{P\in\mathcal E_t}C\bigl(\eta(i(P)),P\bigr)\right]\geq0.$
\end{talkpoints}
\medskip
\begin{equation*}
\boxed{W_t^\mu\geq0\qquad(\mu\in\Mstar).}
\end{equation*}
% Spoken text: patches_talk_proof_script.tex; search for ProofPositivityScript.
\note{\ProofPositivityScript}
\end{frame}

% 19. 15:35--16:30
\begin{frame}[label=proof-split]{Proof of main result}
\begin{equation*}
L_t:=\{\text{no interaction in }(t/2,t]\}\in\mathcal G_t.
\end{equation*}
\bigskip
\begin{equation*}
(\mu P_t)(\chi_A)
=\underbrace{\E_A[W_t^\mu\ind_{L_t}]}_{\text{weak dependence on }\mu}
+\underbrace{\E_A[W_t^\mu\ind_{L_t^c}]}_{\text{small contribution}}.
\end{equation*}
\bigskip
\begin{itemize}
\item The extinction event $\{A_t=\varnothing\}$ is hidden behind the patches.
\medskip
\item $L_t$ is a $\mathcal G_t$-measurable proxy for extinction.
\end{itemize}
% Spoken text: patches_talk_proof_script.tex; search for ProofSplitScript.
\note{\ProofSplitScript}
\end{frame}

% 20. 16:30--18:00
\begin{frame}[label=proof-late]{Late interaction term}
\smallline{Primes denote the system with the independent $1\to0$ transitions at rate $\varepsilon$ removed.}
\medskip
\begin{columns}[T,onlytextwidth]
\begin{column}{.34\textwidth}
\centering
\paperfigure{backward-outgoing-chain}{3.8cm}
\end{column}
\begin{column}{.63\textwidth}
\begin{talkpoints}
\item Same dual process and patches; larger potential.
\item For a patch ending with an outgoing interaction,
\begin{equation*}
C(P)=e^{-\varepsilon\abs P}C'(P).
\end{equation*}
\item For an ancestry reaching $u>t/2$,
\begin{equation*}
\sum_{P\text{ in the ancestry}}\abs P=u,
\qquad e^{-\varepsilon u}\leq e^{-\varepsilon t/2}.
\end{equation*}
\end{talkpoints}
\end{column}
\end{columns}
\medskip
\begin{align*}
0\leq\E_A[W_t^\mu\ind_{L_t^c}]
&\leq e^{-\varepsilon t/2}\E_A[(W_t^\mu)']\\
&=e^{-\varepsilon t/2}(\mu P_t')(\chi_A)
\leq e^{-\varepsilon t/2}.
\end{align*}
% Spoken text: patches_talk_proof_script.tex; search for ProofLateScript.
\note{\ProofLateScript}
\end{frame}

% 21. 18:00--18:45
\begin{frame}[label=proof-end]{No-late-interaction term}
\begin{columns}[T,onlytextwidth]
\begin{column}{.45\textwidth}
\centering
\paperfigure{no-late-end-patches}{3.8cm}
\end{column}
\begin{column}{.52\textwidth}
\begin{talkpoints}
\item On $L_t$, every end patch extends for at least $t/2$.
\item In this event's contribution, dependence on $z$ in $C(z,P)$ decays exponentially with the added length.
\item Errors accumulate over the $\abs{\mathcal E_t}$ end patches.
\end{talkpoints}
\end{column}
\end{columns}
\medskip
\begin{align*}
\abs{\E_A[(W_t^\mu-W_t^\nu)\ind_{L_t}]}
\leq e^{-\varepsilon t/2}\E_A\!\left[
(1+\abs{\mathcal E_t})\bigl((W_t^\mu)'+(W_t^\nu)'\bigr)\ind_{L_t}\right],
\quad \mu,\nu\in\Mstar.
\end{align*}
\smallline{Polynomial graph growth controls the accumulated end-patch error.}
% Spoken text: patches_talk_proof_script.tex; search for ProofEndScript.
\note{\ProofEndScript}
\end{frame}

% 22. 18:45--19:20
\begin{frame}[label=proof-limit]{Proof finish}
\begin{itemize}
\item The two estimates give exponential loss of dependence on the initial law:
\begin{equation*}
\abs{(\mu P_t)(\chi_A)-(\nu P_t)(\chi_A)}\longrightarrow0,
\qquad \mu,\nu\in\Mstar.
\end{equation*}
\medskip
\item The semigroup preserves $\Mstar$, so there is invariant law $\pi\in\Mstar$ that must be the common limit:
\begin{equation*}
\mu P_t\Longrightarrow\pi,
\qquad \mu\in\Mstar.
\end{equation*}
\medskip
\item A linearity argument extends the convergence to every $\mu\in\Mminus$.
\end{itemize}
% Spoken text: patches_talk_proof_script.tex; search for ProofLimitScript.
\note{\ProofLimitScript}
\end{frame}

% 23. 19:20--20:00
\begin{frame}[label=proof-ergodicity]{Upgrade result to ergodicity}
\begin{talkpoints}
\item If $p_i^\star\leq\tfrac12$ for every $i$, then $\mathcal M_{-,1}$ contains every probability law.
\par\medskip
Hence uniform exponential ergodicity.
\item On an orientable neighbourhood graph, a separate burn-in argument gives
\begin{equation*}
\liminf_{t\to\infty}\inf_{\eta\in\Omega}P_t\chi_A^\star(\eta)\geq0
\qquad\text{for every finite }A\subset\Lambda.
\end{equation*}
Every initial law approaches $\Mstar$; hence uniform ergodicity.
\end{talkpoints}
% Spoken text: patches_talk_proof_script.tex; search for ProofErgodicityScript.
\note{\ProofErgodicityScript}
\end{frame}

\appendix
\setbeamertemplate{footline}{\hfill\usebeamerfont{footline}Backup\hspace{3mm}\vspace{2mm}}

\begin{frame}[noframenumbering]{The coefficient criterion}
For $g(\eta)=\sum_S\widehat g(S)\prod_{j\in S}\eta(j)$, patch positivity is
\begin{equation*}
\widehat\lambda_i(S)\leq0,
\qquad
\widehat{\lambda_i(\rho_i-\rho_i(\mathbf0))}(S)\geq0,
\qquad \varnothing\neq S\subseteq N(i).
\end{equation*}
\begin{equation*}
p_i^\star
=\max\!\left(\{0\}\cup
\left\{\frac{\widehat{\lambda_i\rho_i}(S)}{\widehat\lambda_i(S)}:
\varnothing\neq S\subseteq N(i),\ \widehat\lambda_i(S)<0\right\}\right).
\end{equation*}
\medskip
For noisy FA-1f,
\begin{equation*}
\lambda_i=1-\chi_{N(i)}+\varepsilon,\qquad
\lambda_i\rho_i=p(1-\chi_{N(i)})+\varepsilon q.
\end{equation*}
\begin{equation*}
\widehat{\lambda_i(\rho_i-\rho_i(\mathbf0))}\bigl(N(i)\bigr)
=\frac{\varepsilon(q-p)}{1+\varepsilon}.
\end{equation*}
\note{Source: current manuscript, definitions of patch positivity and threshold profile, and the noisy FA-1f example. There is only one nonconstant coefficient of lambda, equal to minus one. The only nonconstant coefficient of lambda rho is minus p. Thus the threshold is p and the second sign condition holds exactly when q is at least p. This calculation does not prove an ergodicity result for q below p.}
\end{frame}

\begin{frame}[noframenumbering]{\texorpdfstring{Why the threshold $\frac12$ reaches every initial law}{Why one half reaches every initial law}}
If $p_i^\star\leq\frac12$, then
\begin{equation*}
\abs{\eta(i)-p_i^\star}\leq1-p_i^\star,
\qquad
\abs{\chi_B^\star(\eta)}\leq\prod_{i\in B}(1-p_i^\star).
\end{equation*}
Hence every probability measure satisfies
\begin{equation*}
\mu(\chi_B^\star)\geq-\prod_{i\in B}(1-p_i^\star),
\qquad\text{so }\mu\in\mathcal M_{-,1}.
\end{equation*}
\bigskip
For a constant threshold $p$, product measures satisfy
\begin{equation*}
\mu_u\in\mathcal M_{-,1}
\quad\Longleftrightarrow\quad u\geq\max\{2p-1,0\}.
\end{equation*}
\note{The one-half threshold comes from the elementary bound on the centered spins, rather than from the geometry of the graph. For the product-measure characterization assume the underlying lattice is infinite, as in the talk.}
\end{frame}

\begin{frame}[noframenumbering]{The signed dual for noisy FA-1f}
\begin{equation*}
\Lcal_\eta H=\Dcal H+VH,
\qquad H((A,\sigma),\eta)=\sigma\chi_A(\eta).
\end{equation*}
At each active $i$, with $S=N(i)$:
\begin{center}
\begin{tabular}{@{}llll@{}}
\toprule
Update & new active set & sign multiplier & rate\\
\midrule
death & $A\setminus\{i\}$ & $+$ & $p+\varepsilon q$\\[1mm]
split & $(A\setminus\{i\})\cup S$ & $-$ & $p$\\[1mm]
birth & $A\cup S$ & $+$ & $1$\\
\bottomrule
\end{tabular}
\end{center}
\begin{equation*}
V(A)=(2p-\varepsilon(1-q))\abs A.
\end{equation*}
\smallline{Removing original $1\to0$ noise changes the potential, while preserving these dual jumps.}
\note{These are the bulk coefficients with the full neighbourhood; boundary coefficients of the zero-boundary restriction may differ. A dual death comes from the constant birth coefficient of the original spin generator. It must not be confused with an original pure-death transition. Removing original deaths leaves lambda rho and all nonconstant lambda coefficients unchanged, and therefore leaves the dual jump rates unchanged. The potential increases.}
\end{frame}

\end{document}



